\subsection{Data and evaluation unit}

We evaluate on the RealNews-based subset used throughout the project. The paragraph is the primary evaluation unit: each source paragraph produces a set of candidate rewrites, and a paragraph counts as a success only if at least one selected candidate satisfies both the detector criterion and the validity criterion. This distinction matters because candidate-level success can be higher than paragraph-level success when multiple candidates are generated for the same source.

The main strict evaluation contains 30 source paragraphs and 120 feedback-driven candidates. Sixteen candidates pass the full strict filter, covering 9 distinct paragraphs. We therefore report both the candidate-level strict rate, 13.3\% (16/120), and the paragraph-level strict-valid ASR, 30.0\% (9/30). Reporting both quantities prevents the candidate pool size from being mistaken for attack success.

\subsection{Detector and decision thresholds}

All comparisons use the same fixed AMR-based detector and the same detector decision threshold, $\tau=2.33$. Lower detector scores indicate stronger evidence of escape, and the detector threshold is only used for the final binary decision. Candidate generation additionally uses a feedback threshold of 0.65 when ranking, retaining, and prioritizing rewrites during search. These are different thresholds with different functions: $\tau$ defines evaluation-time success, while 0.65 controls the search procedure.

To check whether the headline result depends on the search setting, we re-evaluate the stored candidate scores over thresholds from 0.50 to 0.80. The paragraph-level strict-valid ASR remains 30.0\% throughout the 0.65--0.80 interval. This is a sensitivity result on the saved candidates; it does not imply that the same plateau must hold outside the evaluated range or under a different candidate pool.

\subsection{Candidate budgets and comparison systems}

Candidate-budget analysis evaluates $K\in\{1,2,4,8,16\}$. For each $K$, selection is restricted to the first $K$ candidates available for a paragraph and follows the validity-aware rule in Section~3. This separates gains due to the rewriting strategy from gains due only to a larger search budget. The best observed strict-valid ASR occurs at $K=1$ (33.3\%), while larger budgets yield 26.7--30.0\%. In other words, additional search lowers the mean detector score but does not monotonically improve strict validity.

We compare against two existing runs evaluated with the same detector: a 30-prompt pilot baseline and a 250-sample formal baseline. We also include prompt-component ablations---weak paraphrasing, rewriting without boundary changes, boundary-only rewriting, and the full E2 composition---to isolate which structural operations drive detector reduction. These baselines are detector-aligned reference points, but they are not fully matched on sample size, candidate construction, or quality-audit coverage, so they are not controlled head-to-head comparisons.

\subsection{Evaluation measures}

The primary measure is paragraph-level strict-valid ASR. Secondary measures include raw detector escape rate, candidate-level strict rate, mean detector score, mean detector-score reduction, and AUROC where a matched comparison distribution is available. Quality analysis records surface validity, semantic validity, mean review score, the proportion of outputs scoring at least four, and the major-error rate. Repeatability is measured by rerunning detector recomputation 50 times and counting candidates whose binary decisions change.

These measures are reported together because they answer different questions. Raw escape captures detector sensitivity, candidate-level strict rate captures the efficiency of the search process, paragraph-level strict-valid ASR captures end-to-end usefulness, and repeatability captures whether the observed decision is stable under recomputation.
